% Generated draft: requires mathematical and source review.
\documentclass[11pt]{article}
\usepackage[margin=25mm]{geometry}
\usepackage{fontspec}
\setmainfont{Latin Modern Roman}
\usepackage{amsmath,amssymb,mathtools,unicode-math}
\setmathfont{Latin Modern Math}
\usepackage{enumitem,needspace,xurl,hyperref}
\setlength{\parindent}{0pt}
\setlength{\parskip}{5pt}
\setlength{\emergencystretch}{3em}
\newcommand{\N}{\mathbb{N}}
\newcommand{\Z}{\mathbb{Z}}
\newcommand{\Q}{\mathbb{Q}}
\newcommand{\R}{\mathbb{R}}
\newcommand{\C}{\mathbb{C}}
\newcommand{\F}{\mathbb{F}}
\newcommand{\A}{\mathbb{A}}
\newcommand{\PP}{\mathbb{P}}
\newcommand{\E}{\mathbb{E}}
\newcommand{\eps}{\varepsilon}
\newcommand{\dd}{\,\mathrm{d}}
\begin{document}
\section*{UnsolvedMath ID: 30000757}
\textbf{Source problem code:} catalogue entry
\subsection{Definitions and mathematical statement}
Let \(k\) be a number field, \(G/k\) a semiabelian variety, \(V/k\) a smooth variety, and \(\pi:V\to G\) a finite \(k\)-morphism. Semiabelian means an extension of an abelian variety by a torus. Let \(S\) be a finite set of places and
\[
 \mathcal O_S=\{a\in k:v(a)\geq0\text{ for every finite place }v\notin S\}.
\]
The integral-point notation requires a chosen model \(\mathcal V/\mathcal O_S\) with generic fiber \(V\); \(V(\mathcal O_S)\) conventionally denotes the resulting \(\mathcal V(\mathcal O_S)\). The record does not specify this model, and integral points are not simply all \(k\)-points of an arbitrary variety.

The printed conjecture is
\[
 \overline{V(\mathcal O_S)}^{\mathrm{Zar}}=V
 \quad\Longrightarrow\quad
 \pi\text{ is an unramified covering of }G,
\]
and adds that \(V\) should carry a semiabelian structure for which \(\pi\) is an isogeny. An isogeny is a surjective algebraic group homomorphism with finite kernel. The latter statement is stronger than identifying a finite unramified map of varieties; it requires compatible group identities and descent over the specified field.

Two essential qualifications are missing from the supplied hypotheses. A finite morphism is not by definition dominant or surjective, while a covering of \(G\) is. In the customary finite-\(\acute{e}\)tale interpretation, a covering also includes the appropriate flatness and surjectivity requirements, beyond the unramified condition alone. The record states no dominance, dimension or connectedness hypotheses guaranteeing this setup. These cannot be inserted silently in order to obtain the conclusion.

Moreover, for \(\pi\) itself to be an isogeny over \(k\), its fiber over the identity of \(G\) must contain the \(k\)-rational identity of the proposed group structure on \(V\). Density of integral points provides \(k\)-points on \(V\) but does not, as a printed hypothesis, select a rational lift of that identity. A statement after translating the target, choosing basepoints, extending \(k\), or treating \(V\) as a torsor differs from the literal source conclusion. The intended descent/translation convention is therefore explicitly unresolved, not replaced by an automatic ``in particular.''

The literature assessment checked 2026-08-21 says only a declaration of a conjecture survived, whereas both current statements give the full finite-cover assertion. It supplies no assessment resolving the missing dominance or group-structure conventions. The open label and reported complete-source verification remain dated metadata, not independently checked current results.
\subsection{Short English statement}
The source predicts that dense integral points on a smooth finite cover of a semiabelian variety force an unramified covering and an isogeny structure. As printed, dominance and a compatible rational identity are missing, and the integral model is undefined. Those qualifications are retained as gaps rather than silently added.
\subsection{Sources}
\begin{itemize}
\item Local UnsolvedMath registry, permanent ID 30000757, source code catalogue entry, ``Integral Points on Finite Covers of Semiabelian Varieties,'' including the full original/cleaned conjecture and its additional semiabelian-isogeny conclusion.
\item \emph{Diophantische Approximationen (2008)}, Oberwolfach Reports, pp.\ 1115--1190, DOI 10.4171/owr/2007/21: \url{https://doi.org/10.4171/owr/2007/21}. This is the sole supplied literature reference. The citation uses 2008, while the DOI and proposed-year metadata use 2007.
\item Literature triage checked 2026-08-21; record updated 2026-09-16. The assessment describes a missing statement despite the current explicit conjecture. No external source has been visited or independently checked here.
\end{itemize}
\end{document}
